\begin{helper}[HypergeometricCountingTailBound]
\label{helper:HypergeometricCountingTailBound}
Let \(N,n,m\) be natural numbers with \(n\le N\) and \(m\le N\), and let
\(A\subseteq\operatorname{Fin}(N)\) have \(|A|=m\).  Choose an
\(n\)-element subset \(S\subseteq\operatorname{Fin}(N)\) uniformly, and set
\[
Y=|S\cap A|,\qquad q=\frac mN,\qquad \mu=nq,
\]
with the same real-division convention as Lean, so \(q=0\) when \(N=0\).
For every real \(T>0\),
\[
\frac{\left|\{S: |S|=n,\ |S\cap A|\ge T\}\right|}{\binom Nn}
\le \left(\frac{e\mu}{T}\right)^T .
\]
\end{helper}
\begin{proof}
If \(N=0\), then \(n=0\) and \(m=0\).  There is only the empty sample and
\(Y=0\), so \(T>0\) makes the event \(Y\ge T\) empty.  The left side is
\(0\), and the right side is also nonnegative, hence the desired inequality
holds.  We may therefore assume \(N>0\).

If \(\mu=0\), then \(nq=0\).  Since \(N>0\), \(q=m/N\), so either \(n=0\) or
\(m=0\).  In both cases \(Y=0\) for every sample: if \(n=0\) the sample is
empty, and if \(m=0\) the marked set is empty.  Because \(T>0\), the event
\(Y\ge T\) is empty, so the left side is \(0\).  The right side is
\((e\mu/T)^T=0\), and the inequality again holds.

Assume from now on that \(\mu>0\).  The left side is a probability, so it is
at most \(1\).  If \(T\le\mu\), then
\[
\frac{e\mu}{T}\ge e,
\]
and therefore \((e\mu/T)^T\ge1\), since \(T>0\).  The bound follows in this
case.

It remains to handle \(T>\mu>0\).  Set
\[
\lambda=\log(T/\mu).
\]
Then \(\lambda>0\).  Markov's inequality applied to the nonnegative random
variable \(e^{\lambda Y}\) gives
\[
\Pr(Y\ge T)
=\Pr(e^{\lambda Y}\ge e^{\lambda T})
\le \mathbb E[e^{\lambda Y}]\,e^{-\lambda T}.
\]
By \(\noderef{HypergeometricMgfComparison}\), the hypergeometric moment
generating function is bounded by
\[
\mathbb E[e^{\lambda Y}]
\le \exp(\mu(e^\lambda-1)).
\]
Combining the last two displays yields
\[
\Pr(Y\ge T)
\le \exp(\mu(e^\lambda-1)-\lambda T).
\]
The choice of \(\lambda\) gives \(e^\lambda=T/\mu\).  Hence
\[
\mu(e^\lambda-1)-\lambda T
=\mu\left(\frac T\mu-1\right)-T\log(T/\mu)
=T-\mu-T\log(T/\mu).
\]
Since \(\mu\ge0\),
\[
T-\mu-T\log(T/\mu)
\le T-T\log(T/\mu)
=\log\left(\left(\frac{e\mu}{T}\right)^T\right).
\]
Exponentiating both sides gives
\[
\Pr(Y\ge T)\le \left(\frac{e\mu}{T}\right)^T,
\]
which is the required estimate.
\end{proof}
